\documentclass[11pt]{article}

\usepackage{amsmath}
\usepackage{amsthm}
\usepackage{amsfonts}
\usepackage{hyperref}	
\usepackage{xcolor}	
\usepackage{graphicx}
\usepackage{multirow}

\setlength{\textwidth}     {16.0cm}
\setlength{\textheight}    {21.0cm}
\setlength{\evensidemargin}{ 0.0cm}
\setlength{\oddsidemargin} { 0.0cm}
\setlength{\topmargin}     {-0.5cm}
\setlength{\baselineskip}  { 0.7cm}

\newcommand{\R}{\mathbb{R}}
\newcommand{\N}{\mathbb{N}}
\newcommand{\xkjtrial}{x^{k,j}_{\mathrm{trial}}}
\newcommand{\minimize}{\mathop{\text{Minimize}}}
\newcommand{\flow}{f_{\mathrm{low}}}
\newcommand{\simplex}{\textstyle\sum}

\newtheorem{theorem}{Theorem}[section]
\newtheorem{corollary}{Corollary}[section]
\newtheorem{lemma}{Lemma}[section]
\newtheorem{definition}{Definition}[section]
\newtheorem{assumption}{Assumption}
\renewcommand\theassumption{A\arabic{assumption}}

% Show line numbers
\usepackage[mathlines]{lineno}
\renewcommand\linenumberfont{\normalfont\small\color{gray}}
\linenumbers

% Show labels
\usepackage{showlabels}


\begin{document}

\title{A quasi-Newton regularized method for box-constrained order-value optimization\thanks{This work has been funded by the Vice-Rectorate for Research, Outreach and Social Projection of the University of the Atlantic through the 2022 Internal Call for Proposals to Strengthen Research Groups with Tools for Creation, Innovation, and Research.}}

\author{
    G. Q. Álvarez\thanks{Mathematics Program, Faculty of Basic Sciences, University of Atlantic, Puerto Colombia, Atlántico, Colombia. e-mail: gdquintero@uniatlantico.edu.co}
    \and 
    Carlos Araujo\thanks{Mathematics Program, Faculty of Basic Sciences, University of Atlantic, Puerto Colombia, Atlántico, Colombia. e-mail: carlosaraujo@uniatlantico.edu.co}
    \and
    M. A. C. Candezano\thanks{Mathematics Program, Faculty of Basic Sciences, University of Atlantic, Puerto Colombia, Atlántico, Colombia. e-mail: miguelcaro@uniatlantico.edu.co}
}

% \date{June 18, 2025\footnote{Revision made on \today.}}
\date{\today}
\maketitle

\begin{abstract}
The order-value function of order $p$ returns, at each point, the $p$th smallest
value among a finite family of functions $f_1,\dots,f_m$; minimizing it recovers
the minimization of the minimum and of the maximum of the family in the extreme
cases $p=1$ and $p=m$, and provides a flexible tool for robust estimation, in
which a prescribed number of outlying observations are discarded. We propose a
regularized method for minimizing the order-value function subject to box
constraints. At each iteration the method minimizes a strongly convex model that
combines a piecewise-linear term built from the active gradients, a quasi-Newton
curvature term defined by a symmetric positive semidefinite matrix, and a
quadratic regularization; the curvature term carries second-order information
without requiring explicit second derivatives, and a purely first-order
regularized method is recovered as the special case in which this term vanishes.
We introduce an approximate optimality condition for the problem and prove that
the method is well defined and attains the same worst-case iteration- and
evaluation-complexity bounds as its first-order counterpart, so that the
incorporation of curvature preserves the theoretical guarantees at no additional
order of complexity. The applicability of the method to parameter-estimation
problems with outliers is illustrated, and the practical effect of the curvature
is examined empirically.\\

\noindent
\textbf{Keywords:} Order-value optimization, regularized models,
quasi-Newton methods, complexity, algorithms, applications.\\

\noindent
\textbf{Mathematics Subject Classification:} 90C30, 65K05, 49M37, 90C60, 68Q25.\\
\end{abstract}

\section{Introduction} \label{sec:intro}

Many optimization problems involve objective functions that are obtained not
by aggregating a collection of simpler functions, but by selecting one of them
according to its rank. The prototypical example is the order-value function of
order $p$, introduced in \cite{dunder1,dunder2}: given $m$ functions
$f_1,\dots,f_m$, its value at a point $x$ is the $p$th smallest among
$f_1(x),\dots,f_m(x)$. Functions of this kind belong to the broader class of
generalized order-value functions---those whose value at $x$ is governed by
order relations among the elements of a set $\{f_i(x)\}_{i\in I}$---a class
systematized in \cite{mtop}. Since the index that attains the $p$th position
may switch abruptly as $x$ moves, the order-value function is generally
nonsmooth even when each $f_i$ is continuously differentiable.

The extreme choices $p=1$ and $p=m$ recover the minimum and the maximum of
$\{f_i(x)\}$, respectively, so the order-value function interpolates between
these two classical aggregates; the intermediate cases are the ones of
practical interest. When $x$ encodes portfolio positions and $f_i(x)$ is the
loss incurred under scenario $i$, the order-value function coincides with the
discrete value-at-risk (VaR) measure, a standard risk measure in financial
risk management \cite{jorion}. When $f_i(x)$ measures the misfit of a
parametric model to the $i$th observation, minimizing the $p$th order-value
function fits the parameters while discarding the $m-p$ worst-fitted
observations. This makes order-value optimization a natural tool for robust
estimation, where one seeks to remain insensitive to a minority of grossly
corrupted data, or outliers.

The order-value optimization (OVO) problem was first formulated in
\cite{dunder1}, where a steepest-descent-type method was proposed and shown to
converge to points satisfying a weak optimality condition. Sharper optimality
conditions and a reformulation as a nonlinear program with equilibrium
constraints were subsequently developed in \cite{dunder2}. A quasi-Newton
method generalizing the Cauchy scheme of \cite{dunder1}---approximating the
active components by convex quadratics and proving global and local
convergence---was introduced in \cite{yano1}, and a global optimization
strategy combining multistart with a tunneling technique was proposed in
\cite{salva}. A closely related variant, in which the sum of the $p$ smallest
values is minimized---low order-value optimization (LOVO)---was developed with
applications to robust estimation and data analysis in \cite{yano2}. On the
computational-complexity side, \cite{jiang} established that minimizing the
order-value function over a polytope is NP-hard in the strong sense.

The systematic study of worst-case complexity in continuous optimization
originates in the cubic-regularization analysis of \cite{nesterov}, and
complexity guarantees have since been obtained for a broad range of problems
(see, e.g., \cite{bgmst2,bgmst1} and the monograph \cite{cogtbook}). Within
this line, \cite{alvarezovo} recently proposed a first-order method for
box-constrained OVO that minimizes, at each iteration, a quadratically
regularized piecewise-linear model of $f$, and established iteration- and
evaluation-complexity bounds for reaching an approximate optimality condition.

In the present work we generalize the approach of \cite{alvarezovo} by
enriching the regularized model with second-order information. At each
iteration our method minimizes a model built from the same piecewise-linear
convex term, a curvature term governed by a symmetric positive semidefinite
matrix $B$, and a Tikhonov-type quadratic regularization. The matrix $B$ is
updated in a quasi-Newton fashion and therefore requires no explicit second
derivatives, while the regularization renders the model strongly convex, so
that each subproblem has a unique minimizer. Setting $B=0$ recovers the
first-order method of \cite{alvarezovo} exactly, so the proposed method is a
strict generalization of it. Curvature had already been incorporated into OVO
by the quasi-Newton method of \cite{yano1}, which approximates the active
components by convex quadratics with positive semidefinite quasi-Newton
Hessians and establishes global convergence together with local superlinear
and quadratic rates; that method, however, controls the step through a
trust-region bound combined with a line search, and its guarantees are
asymptotic. Here the curvature is instead embedded in the adaptively
regularized model of \cite{alvarezovo}: the Tikhonov term renders each
subproblem strongly convex with a unique minimizer, and the analysis yields
non-asymptotic worst-case iteration- and evaluation-complexity bounds rather
than asymptotic convergence rates. We derive an approximate optimality
condition for
the box-constrained problem and prove that the method is well defined and
attains the \emph{same} worst-case iteration- and evaluation-complexity bounds
as the first-order method. Thus the incorporation of curvature does not degrade
the theoretical guarantees; whether it also yields a practical advantage is a
separate question, which we examine empirically on parameter-estimation
problems in the presence of outliers.

The remainder of the paper is organized as follows. Section~\ref{sec:method}
defines the box-constrained order-value optimization problem and the proposed
regularized method. Section~\ref{sec:complexity} introduces an appropriate
optimality condition and establishes the well-definiteness, convergence, and
complexity of the method. Section~\ref{sec:numerical} reports illustrative
numerical experiments on parameter estimation in the presence of outliers.
Final remarks are given in Section~\ref{sec:final}.

\medskip
\noindent\textit{Notation.} The symbol $\|\cdot\|$ denotes the Euclidean norm.
For $i=1,\dots,n$, $e_i\in\mathbb{R}^{n}$ denotes the $i$th column of the
identity matrix.

\section{A quasi-Newton regularized method} \label{sec:method}

Let $f_i : \mathbb{R}^n \to \mathbb{R}$, $i = 1, \ldots, m$, be given and let
$p \in \{1, \ldots, m\}$ be fixed. The \textit{order-value function of order
$p$} is the mapping $f : \mathbb{R}^n \to \mathbb{R}$ defined by
\begin{equation}
    f(x) = f_{i_p(x)}(x), \label{ovo-function}
\end{equation}
where $\{i_1(x), \ldots, i_m(x)\}$ is a permutation of $\{1, \ldots, m\}$ such
that
\begin{equation}
    f_{i_1(x)}(x) \leq f_{i_2(x)}(x) \leq \cdots \leq f_{i_m(x)}(x);
    \label{fi-sort}
\end{equation}
that is, $f(x)$ is the $p$th smallest value among $f_1(x), \ldots, f_m(x)$.
The boundary cases $p = 1$ and $p = m$ recover, respectively, the minimum and
the maximum of $\{f_1(x), \ldots, f_m(x)\}$.

If $f_1, \ldots, f_m$ are continuous, then so is $f$. Differentiability,
however, is not inherited: $f$ may be nonsmooth at a point $x$ at which the
$p$th smallest value is attained by more than one index, i.e.\ whenever there
exist $i \neq j$ with $f_i(x) = f_j(x) = f(x)$, since the active index $i_p(x)$
may then change as $x$ varies.

We address the box-constrained order-value optimization problem
\begin{equation}
    \minimize_{x \in \Omega} f(x), \label{ovo-problem}
\end{equation}
where $\Omega := \{x \in \mathbb{R}^n \mid \ell \leq x \leq u\}$ with
$\ell, u \in \mathbb{R}^n$ and $\ell_i < u_i$ for $i = 1, \ldots, n$. Box
constraints arise naturally in parameter estimation, where the unknowns are
confined to physically or statistically meaningful bounds.

For $\delta > 0$ and $x \in \Omega$, define the \textit{active index set}
\begin{equation}
    I_\delta(x) := \{i \in \{1, \ldots, m\} \mid
    f(x) - \delta \leq f_i(x) \leq f(x) + \delta\}, \label{idelta}
\end{equation}
the set of indices whose values lie within a band of half-width $\delta$ around
$f(x)$. As $\delta \to 0$, $I_\delta(x)$ shrinks to the set of indices that
attain the order-value,
\begin{equation*}
    I_0(x) := \{i \in \{1, \ldots, m\} \mid f_i(x) = f(x)\}.
\end{equation*}

The method builds, around a reference point $\bar{x} \in \Omega$, a local model
of $f$ that combines the first-order information of the active components with
curvature carried by a symmetric positive semidefinite matrix
$B \in \mathbb{R}^{n \times n}$. In a quasi-Newton framework, $B$ is updated
from successive gradients, so that no second derivatives are computed
explicitly. For given $\bar{x} \in \Omega$, $B \succeq 0$, $\delta > 0$, and
$\sigma > 0$, the regularized model is
\begin{equation}
    \Psi(x; \bar{x}, B, \delta, \sigma) =
    \max_{i \in I_\delta(\bar{x})}
    \left\lbrace \nabla f_i(\bar{x})^T(x - \bar{x}) 
    + \frac{1}{2} \left[ (x - \bar{x})^T B (x - \bar{x})
    + \sigma \|x - \bar{x}\|^2 \right]\right\rbrace. \label{model}
\end{equation}
The first term is the piecewise-linear convex function that captures the
worst-case first-order behavior of $f$ over the active components; the bracket
adds the quasi-Newton curvature $(x - \bar{x})^T B (x - \bar{x})$ and a
Tikhonov regularization $\sigma \|x - \bar{x}\|^2$. Since $\sigma > 0$, the
model is strongly convex even when $B$ is only positive semidefinite, so the
subproblem \eqref{ovo-subpro} has a unique minimizer. For $B = 0$, the model
reduces to the first-order regularized model of \cite{alvarezovo}, which the
present method thus generalizes.

\medskip
\noindent\textbf{Algorithm~\ref{sec:method}.1.} Let $\delta > 0$,
$\sigma_{\min} > 0$, $M > 0$, $\alpha \in (0,1)$, $\gamma > 1$, and
$x^0 \in \Omega$ be given. Initialize $k \leftarrow 0$.
\begin{description}
\item[Step~1.] Initialize $j \leftarrow 0$ and $\sigma_{k,j} = \sigma_{\min}$.
\item[Step~2.] Choose a symmetric positive semidefinite matrix
$B_{k,j} \in \mathbb{R}^{n \times n}$ with $\|B_{k,j}\| \leq M$ and compute
$\xkjtrial$ as the global minimizer of
    \begin{equation} \label{ovo-subpro}
    \minimize_{x \in \Omega}\, \Psi(x; x^k, B_{k,j}, \delta, \sigma_{k,j}).
    \end{equation}
\item[Step~3.] If
    \begin{equation} \label{ovo-armijo}
    f(\xkjtrial) \leq f(x^k) - \alpha \|\xkjtrial - x^k\|^2
    \end{equation}
does not hold, set $\sigma_{k,j+1} = \gamma \sigma_{k,j}$,
$j \leftarrow j + 1$, and go to Step~2.
\item[Step~4.] Set $x^{k+1} = \xkjtrial$, $\sigma_k = \sigma_{k,j}$,
$j_k = j$, $k \leftarrow k+1$, and go to Step~1.
\end{description}

The bound $\|B_{k,j}\| \leq M$ keeps the curvature from dominating the model.
The freedom in choosing $B_{k,j}$ lets the method accommodate standard
quasi-Newton updates such as BFGS or SR1, and the choice $B_{k,j} = 0$ for all
$k, j$ recovers the first-order method of \cite{alvarezovo}.

The requirements on $\xkjtrial$ are mild. The theory only needs $\xkjtrial$ to
be a stationary point of \eqref{ovo-subpro} with
$\Psi(\xkjtrial; x^k, B_{k,j}, \delta, \sigma_{k,j}) \leq 0$; since
$\Psi(x^k; x^k, B_{k,j}, \delta, \sigma_{k,j}) = 0$, any descent method started
at $x^k$ produces such iterates. Moreover, \eqref{ovo-subpro} minimizes a
strongly convex function over the compact convex set $\Omega$, so it has a
unique global minimizer and every stationary point is global; the
global-minimizer requirement of Step~2 is therefore met by any such descent
method, and standard convex solvers apply directly.

\section{Optimality condition and complexity} \label{sec:complexity}

\begin{assumption} \label{a1}
The functions $f_1, \ldots, f_m$ belong to $\mathcal{C}^1(\Omega)$; that is,
they are continuously differentiable on $\Omega$.
\end{assumption}

We first show that a local minimizer of \eqref{ovo-problem} also minimizes the
model $\Psi$ built at the corresponding reference point. These results connect
the original nonsmooth problem with the convex subproblems solved at each
iteration and underlie the optimality condition introduced below.

\begin{theorem} \label{teo:conv1}
Suppose that Assumption~\ref{a1} holds. If $x^*$ is a local minimizer of
\eqref{ovo-problem}, then it is also a local minimizer of
$\minimize_{x \in \Omega}\, \Psi(x; x^*, 0, 0, 0)$.
\end{theorem}

\begin{proof}
Suppose not. Then there is $y^* \in \Omega$ with
$\Psi(y^*; x^*, 0, 0, 0) < \Psi(x^*; x^*, 0, 0, 0) = 0$, i.e.\
$\nabla f_i(x^*)^T(y^* - x^*) < 0$ for all $i \in I_0(x^*)$. By
Assumption~\ref{a1}, for each such $i$,
\[
\lim_{t \to 0^+} \frac{f_i(x^* + t(y^* - x^*)) - f_i(x^*)}{t}
= \nabla f_i(x^*)^T(y^* - x^*) < 0,
\]
so there is $\bar{t}_i > 0$ with $f_i(x^* + t(y^* - x^*)) < f_i(x^*) = f(x^*)$
for $t \in (0, \bar{t}_i]$. With $\bar{t} = \min_{i \in I_0(x^*)} \bar{t}_i$,
\begin{equation} \label{bajaste}
f_i(x^* + t(y^* - x^*)) < f(x^*) \quad
\text{for all } i \in I_0(x^*),\ t \in (0, \bar{t}].
\end{equation}
For $j \notin I_0(x^*)$, continuity of $f_1, \ldots, f_m$ gives, for $t$ small,
\begin{align}
f_i(x^* + t(y^* - x^*)) &< f_j(x^* + t(y^* - x^*))
&&\text{if } i \in I_0(x^*),\ f(x^*) < f_j(x^*), \label{arriba}\\
f_i(x^* + t(y^* - x^*)) &> f_j(x^* + t(y^* - x^*))
&&\text{if } i \in I_0(x^*),\ f(x^*) > f_j(x^*). \label{abajo}
\end{align}
Hence the ranking positions of the active indices $I_0(x^*)$ within
$\{f_i(x^* + t(y^* - x^*))\}_{i=1}^m$ are preserved for $t$ small, so the $p$th
position is still occupied by some $i \in I_0(x^*)$ and
$f(x^* + t(y^* - x^*)) = f_i(x^* + t(y^* - x^*))$. With \eqref{bajaste},
$f(x^* + t(y^* - x^*)) < f(x^*)$ for $t$ small. Since $\Omega$ is convex,
$x^* + t(y^* - x^*) \in \Omega$, contradicting the local minimality of $x^*$.
\end{proof}

\begin{theorem} \label{teo:conv2}
Suppose that Assumption~\ref{a1} holds and let $B \succeq 0$ and $\sigma \geq 0$.
If $x^*$ is a local minimizer of \eqref{ovo-problem}, then it is also a local
minimizer of $\minimize_{x \in \Omega}\, \Psi(x; x^*, B, 0, \sigma)$.
\end{theorem}

\begin{proof}
Suppose not: there is $y^* \in \Omega$ with $\Psi(y^*; x^*, B, 0, \sigma) < 0$,
i.e.
\[
\max_{i \in I_0(x^*)} \{ \nabla f_i(x^*)^T(y^* - x^*) \}
+ \tfrac{1}{2} \big[ (y^* - x^*)^T B (y^* - x^*) + \sigma \|y^* - x^*\|^2 \big]
< 0.
\]
Since $B \succeq 0$ and $\sigma \geq 0$, the bracket is nonnegative, so
$\max_{i \in I_0(x^*)} \{ \nabla f_i(x^*)^T(y^* - x^*) \} < 0$, that is,
$\Psi(y^*; x^*, 0, 0, 0) < 0$, contradicting Theorem~\ref{teo:conv1}.
\end{proof}

\begin{theorem} \label{teo:conv3}
Suppose that Assumption~\ref{a1} holds and let $B \succeq 0$ and
$\sigma, \delta \geq 0$. If $x^*$ is a local minimizer of \eqref{ovo-problem},
then it is also a local minimizer of
$\minimize_{x \in \Omega}\, \Psi(x; x^*, B, \delta, \sigma)$.
\end{theorem}

\begin{proof}
Suppose not: there is $y^* \in \Omega$ with
$\Psi(y^*; x^*, B, \delta, \sigma) < 0$. Since
$I_0(x^*) \subseteq I_\delta(x^*)$, we have
$\max_{i \in I_0(x^*)}\{\cdot\} \leq \max_{i \in I_\delta(x^*)}\{\cdot\}$, hence
\[
\max_{i \in I_0(x^*)} \{ \nabla f_i(x^*)^T(y^* - x^*) \}
+ \tfrac{1}{2} \big[ (y^* - x^*)^T B (y^* - x^*) + \sigma \|y^* - x^*\|^2 \big]
\leq \Psi(y^*; x^*, B, \delta, \sigma) < 0,
\]
that is, $\Psi(y^*; x^*, B, 0, \sigma) < 0$, contradicting
Theorem~\ref{teo:conv2}.
\end{proof}

Throughout, $\Sigma_q := \{ \lambda \in \mathbb{R}^q \mid
\sum_{j=1}^q \lambda_j = 1,\ \lambda_j \geq 0 \}$ denotes the unit simplex.
Following \cite{alvarezovo}, the next consequence of Theorem~\ref{teo:conv3}
provides a first-order necessary condition for \eqref{ovo-problem} and
motivates the stopping criterion below.

\begin{corollary} \label{cor:exact}
Suppose that Assumption~\ref{a1} holds and let $\delta \geq 0$. If $x^*$ is a
local minimizer of \eqref{ovo-problem}, then there exist
$\mu \in \Sigma_{|I_\delta(x^*)|}$ and $\nu^\ell, \nu^u \in \mathbb{R}^n_+$ such
that $(\ell_i - x^*_i)\nu^\ell_i = (x^*_i - u_i)\nu^u_i = 0$ for
$i = 1, \ldots, n$ and
\begin{equation} \label{exact-opt}
\sum_{i \in I_\delta(x^*)} \mu_i \nabla f_i(x^*)
+ \sum_{i=1}^n (\nu^u_i - \nu^\ell_i) e_i = 0.
\end{equation}
\end{corollary}

\begin{proof}
By Theorem~\ref{teo:conv3} with $B = 0$ and $\sigma = 0$, $x^*$ minimizes
$x \mapsto \max_{i \in I_\delta(x^*)} \{ \nabla f_i(x^*)^T(x - x^*) \}$ over
$\Omega$. Applying the KKT conditions for the minimization of a maximum of
differentiable functions over a box \cite[Lemma~3.1]{alvarezovo} yields
multipliers $\mu \in \Sigma_{|I_\delta(x^*)|}$ and
$\nu^\ell, \nu^u \in \mathbb{R}^n_+$ with the stated complementarity and
\eqref{exact-opt}.
\end{proof}

\begin{definition} \label{def:optcond}
Given $\epsilon > 0$ and $\delta > 0$, we say that $x \in \Omega$ satisfies the
\textit{approximate optimality condition} $C_{\delta,\epsilon}$ if there exist
$\mu \in \Sigma_{|I_\delta(x)|}$ and $\nu^\ell, \nu^u \in \mathbb{R}^n_+$ such
that $(\ell_i - x_i)\nu^\ell_i = (x_i - u_i)\nu^u_i = 0$ for $i = 1, \ldots, n$
and
\begin{equation} \label{miaprox-box}
\left\| \sum_{i \in I_\delta(x)} \mu_i \nabla f_i(x)
+ \sum_{i=1}^{n} (\nu^u_i - \nu^\ell_i) e_i \right\| \leq \epsilon.
\end{equation}
\end{definition}

Condition $C_{\delta,\epsilon}$ requires the gradient combination of the active
components, corrected by the box multipliers, to be small in norm. For
$\epsilon = 0$ it reduces to the exact condition \eqref{exact-opt} of
Corollary~\ref{cor:exact}. It serves as the stopping criterion for
Algorithm~\ref{sec:method}.1.

\begin{corollary} \label{cor:step}
Suppose that Assumption~\ref{a1} holds. Then, for every $k$ and
$j = 0, \ldots, j_k$, there exist $\mu \in \Sigma_{|I_\delta(x^k)|}$ and
$\nu^\ell, \nu^u \in \mathbb{R}^n_+$ such that
$(\ell_i - [\xkjtrial]_i)\nu^\ell_i = ([\xkjtrial]_i - u_i)\nu^u_i = 0$ for
$i = 1, \ldots, n$ and
\begin{equation} \label{eqcoro1}
\xkjtrial - x^k = \big( B_{k,j} + \sigma_{k,j} I \big)^{-1}
\left[ \sum_{i=1}^{n} (\nu^\ell_i - \nu^u_i) e_i
- \sum_{i \in I_\delta(x^k)} \mu_i \nabla f_i(x^k) \right].
\end{equation}
\end{corollary}

\begin{proof}
Subproblem \eqref{ovo-subpro} minimizes the strongly convex function
$\Psi(\cdot\,; x^k, B_{k,j}, \delta, \sigma_{k,j})$ over the box $\Omega$. By
the KKT conditions for a maximum of differentiable functions over a box
\cite[Lemma~3.1]{alvarezovo}, there exist $\mu \in \Sigma_{|I_\delta(x^k)|}$ and
$\nu^\ell, \nu^u \in \mathbb{R}^n_+$ with the stated complementarity such that
\[
\sum_{i \in I_\delta(x^k)} \mu_i
\big[ \nabla f_i(x^k) + (B_{k,j} + \sigma_{k,j} I)(\xkjtrial - x^k) \big]
+ \sum_{i=1}^n (\nu^u_i - \nu^\ell_i) e_i = 0,
\]
where $\nabla f_i(x^k) + (B_{k,j} + \sigma_{k,j} I)(\xkjtrial - x^k)$ is the
gradient at $\xkjtrial$ of the $i$th component of $\Psi$. Since
$\sum_i \mu_i = 1$, the curvature term factors out; as $\sigma_{k,j} > 0$ and
$B_{k,j} \succeq 0$, the matrix $B_{k,j} + \sigma_{k,j} I$ is positive definite,
and solving for $\xkjtrial - x^k$ gives \eqref{eqcoro1}.
\end{proof}

\begin{assumption} \label{a2}
For all $k$ and $j = 0, \ldots, j_k$, the multipliers $\nu^\ell, \nu^u$ of
Corollary~\ref{cor:step} satisfy
$\max_{i = 1, \ldots, n} |\nu^\ell_i - \nu^u_i| \leq c_\nu$ for some constant
$c_\nu > 0$ independent of $k$ and $j$.
\end{assumption}

The Mangasarian--Fromovitz constraint qualification guarantees that, for each
$(k,j)$, these multipliers are bounded \cite{gauvin}; Assumption~\ref{a2}
strengthens this to a single bound uniform in $k$ and $j$, which is natural
since $\Omega$ is compact and the active gradients are bounded under
Assumption~\ref{a1}.

\begin{assumption} \label{a3}
There is a constant $c_\nabla > 0$ such that $\|\nabla f_i(x)\| \leq c_\nabla$
for all $i = 1, \ldots, m$ and all $x \in \Omega$.
\end{assumption}

Assumption~\ref{a3} holds automatically under Assumption~\ref{a1}, since
$\Omega$ is bounded. Combining Corollary~\ref{cor:step} with
Assumptions~\ref{a2} and~\ref{a3} bounds the step. Since $B_{k,j} \succeq 0$,
the eigenvalues of $B_{k,j} + \sigma_{k,j} I$ are at least $\sigma_{k,j}$, so
$\|(B_{k,j} + \sigma_{k,j} I)^{-1}\| \leq 1/\sigma_{k,j}$, and \eqref{eqcoro1}
gives
\[
\|\xkjtrial - x^k\| \leq \frac{1}{\sigma_{k,j}}
\left( \sum_{i=1}^n |\nu^\ell_i - \nu^u_i|
+ \sum_{i \in I_\delta(x^k)} \mu_i \|\nabla f_i(x^k)\| \right)
\leq \frac{n\,c_\nu + c_\nabla}{\sigma_{k,j}},
\]
where the last inequality uses Assumptions~\ref{a2}, \ref{a3} and
$\mu \in \Sigma_{|I_\delta(x^k)|}$. Thus
\begin{equation} \label{ccotasub}
\|\xkjtrial - x^k\| \leq \frac{c_x}{\sigma_{k,j}}, \qquad
c_x := n\,c_\nu + c_\nabla.
\end{equation}
The step shrinks as $\sigma_{k,j}$ grows, which is the mechanism driving the
acceptance of \eqref{ovo-armijo}.

\begin{assumption} \label{a4}
There is $L > 0$ such that $\|\nabla f_i(y) - \nabla f_i(x)\| \leq L \|y - x\|$
for all $i = 1, \ldots, m$ and $x, y \in \Omega$.
\end{assumption}

Assumption~\ref{a4} is the standard gradient-Lipschitz condition. It implies,
for all $i$ and $x, y \in \Omega$,
\begin{equation} \label{taylor}
f_i(y) \leq f_i(x) + \nabla f_i(x)^T(y - x) + \tfrac{L}{2}\|y - x\|^2,
\end{equation}
as well as the two-sided bound
$|f_i(y) - f_i(x) - \nabla f_i(x)^T(y - x)| \leq \tfrac{L}{2}\|y - x\|^2$. These
make $\Psi$ an upper model of $f$ once $\sigma_{k,j}$ is large relative to $L$,
which yields finite termination of the inner loop. 

\begin{lemma} \label{lem1}
Let $a_1, \ldots, a_r, b_1, \ldots, b_r \in \mathbb{R}$ and $\beta > 0$ satisfy
$a_1 \leq \cdots \leq a_r$, $b_j \leq a_j - \beta$ for $j = 1, \ldots, r$, and
$b_{i_1} \leq \cdots \leq b_{i_r}$ for a permutation $\{i_1, \ldots, i_r\}$ of
$\{1, \ldots, r\}$. Then $b_{i_q} \leq a_q - \beta$ for $q = 1, \ldots, r$.
\end{lemma}

\begin{proof}
See \cite[Lemma~2.1]{dunder1}.
\end{proof}

The next result gives a threshold on $\sigma_{k,j}$ beyond which the trial point
is accepted at Step~3, so that the inner loop terminates after finitely many
increases of $\sigma_{k,j}$, independently of $k$. The threshold scales with
$\delta^{-1}$: a narrower active band demands stronger regularization.

\begin{theorem} \label{teo:suff-descent}
Suppose that Assumptions~\ref{a1}, \ref{a3}, and \ref{a4} hold. Then, for all
$k$ and $j = 0, \ldots, j_k$, if
\begin{equation} \label{este}
\sigma_{k,j} \geq \max\{ L + 2\alpha,\ 9c_x^2/(2\delta) \},
\end{equation}
then $\xkjtrial$ satisfies \eqref{ovo-armijo}.
\end{theorem}

\begin{proof}
Assume \eqref{este}; in particular $\sigma_{k,j} \geq L + 2\alpha$. By
\eqref{taylor}, for every $i \in \{1, \ldots, m\}$,
\begin{align*}
f_i(\xkjtrial)
&\leq f_i(x^k) + \nabla f_i(x^k)^T(\xkjtrial - x^k)
+ \tfrac{L}{2}\|\xkjtrial - x^k\|^2 \\
&\leq f_i(x^k) + \nabla f_i(x^k)^T(\xkjtrial - x^k)
+ \Big( \tfrac{\sigma_{k,j}}{2} - \alpha \Big) \|\xkjtrial - x^k\|^2 \\
&= f_i(x^k) + \nabla f_i(x^k)^T(\xkjtrial - x^k)
+ \tfrac{1}{2} \big[ (\xkjtrial - x^k)^T B_{k,j} (\xkjtrial - x^k)
+ \sigma_{k,j}\|\xkjtrial - x^k\|^2 \big] \\
&\quad - \alpha \|\xkjtrial - x^k\|^2
- \tfrac{1}{2}(\xkjtrial - x^k)^T B_{k,j} (\xkjtrial - x^k).
\end{align*}
Since $B_{k,j} \succeq 0$, the last term is nonpositive and may be dropped. For
$i \in I_\delta(x^k)$, the minimizer property
$\Psi(\xkjtrial; x^k, B_{k,j}, \delta, \sigma_{k,j}) \leq
\Psi(x^k; x^k, B_{k,j}, \delta, \sigma_{k,j}) = 0$ gives
\[
\nabla f_i(x^k)^T(\xkjtrial - x^k)
+ \tfrac{1}{2} \big[ (\xkjtrial - x^k)^T B_{k,j} (\xkjtrial - x^k)
+ \sigma_{k,j}\|\xkjtrial - x^k\|^2 \big] \leq 0,
\]
so that
\begin{equation} \label{cota-lem2}
f_i(\xkjtrial) \leq f_i(x^k) - \alpha \|\xkjtrial - x^k\|^2, \qquad
i \in I_\delta(x^k).
\end{equation}
On the other hand, by \eqref{ccotasub}, the Cauchy--Schwarz inequality, and the
two-sided consequence of Assumption~\ref{a4}, using
$\|\nabla f_i(x^k)\| \leq c_\nabla \leq c_x$ and \eqref{ccotasub},
$|f_i(\xkjtrial) - f_i(x^k)| \leq c_x^2/\sigma_{k,j}
+ L c_x^2/(2\sigma_{k,j}^2)$ for all $i$; using $L \leq \sigma_{k,j}$, this is
at most $3c_x^2/(2\sigma_{k,j})$, and \eqref{este} yields
\begin{equation} \label{teo7eq1}
|f_i(\xkjtrial) - f_i(x^k)| \leq \tfrac{\delta}{3}, \qquad i = 1, \ldots, m.
\end{equation}
Consequently, any index with $f_i(x^k) < f(x^k) - \delta$ satisfies
$f_i(\xkjtrial) < f(x^k)$, and any index with $f_i(x^k) > f(x^k) + \delta$
satisfies $f_i(\xkjtrial) > f(x^k)$. In particular $i_p(\xkjtrial) \in
I_\delta(x^k)$, and the number of components lying below the band is the same at
$x^k$ and $\xkjtrial$. Write $I_\delta(x^k) = \{i_1, \ldots, i_r\}$ with
$f_{i_1}(x^k) \leq \cdots \leq f_{i_r}(x^k)$, and let
$\{i'_1, \ldots, i'_r\}$ reorder them so that
$f_{i'_1}(\xkjtrial) \leq \cdots \leq f_{i'_r}(\xkjtrial)$. Since the count below
the band is preserved and $i_p(\cdot)$ lies in the band at both points,
$i_p(x^k)$ and $i_p(\xkjtrial)$ occupy the same position $q$ within the band;
that is, $i_p(x^k) = i_q$ and $i_p(\xkjtrial) = i'_q$. Applying
Lemma~\ref{lem1} with $a_j = f_{i_j}(x^k)$, $b_j = f_{i_j}(\xkjtrial)$, and
$\beta = \alpha\|\xkjtrial - x^k\|^2$ ---valid by \eqref{cota-lem2}--- gives
$f_{i'_q}(\xkjtrial) \leq f_{i_q}(x^k) - \beta$. Since
$f_{i_q}(x^k) = f_{i_p(x^k)}(x^k) = f(x^k)$,
\[
f(\xkjtrial) = f_{i_p(\xkjtrial)}(\xkjtrial) = f_{i'_q}(\xkjtrial)
\leq f(x^k) - \alpha \|\xkjtrial - x^k\|^2,
\]
which is \eqref{ovo-armijo}.
\end{proof}

The inner loop thus terminates finitely at every outer iteration. The curvature
affects only the constants below ---through the bound $M$ on $\|B_{k,j}\|$--- and
not the $O(\delta^{-2}\epsilon^{-2})$ iteration order, which coincides with that
of the first-order method \cite{alvarezovo}.

\begin{theorem} \label{teo:complexity}
Suppose that Assumptions~\ref{a1}, \ref{a3}, and \ref{a4} hold and that there is
$\flow \in \mathbb{R}$ with $f(x) \geq \flow$ for all $x \in \Omega$. Then
\begin{equation} \label{cota-sigma}
\sigma_k \leq \gamma \max\{ L + 2\alpha,\ 9c_x^2/(2\delta) \},
\end{equation}
and at most
\begin{equation} \label{cota-fcnt}
\big\lfloor 1 + \log_\gamma( \sigma_k / \sigma_{\min} ) \big\rfloor
\end{equation}
function evaluations are performed per outer iteration to obtain
\eqref{ovo-armijo}. Moreover, the number of outer iterations $k$ at which
$C_{\delta,\epsilon}$ is not satisfied by $x^{k+1}$ is at most
\begin{equation} \label{cota-iter}
\left\lfloor \frac{\big( M + \gamma \max\{ L + 2\alpha,\ 9c_x^2/(2\delta) \}
\big)^2}{\alpha} \cdot \frac{f(x^0) - \flow}{\epsilon^2} \right\rfloor.
\end{equation}
\end{theorem}

\begin{proof}
Bounds \eqref{cota-sigma} and \eqref{cota-fcnt} follow from
Theorem~\ref{teo:suff-descent} and the rule
$\sigma_{k,j+1} = \gamma\sigma_{k,j}$: starting from
$\sigma_{k,0} = \sigma_{\min}$, at most
$\lfloor 1 + \log_\gamma(\sigma_k/\sigma_{\min}) \rfloor$ increases occur before
\eqref{este} holds, after which \eqref{ovo-armijo} is satisfied.

Let $K$ be the set of outer indices $k$ at which $C_{\delta,\epsilon}$ fails at
$x^{k+1}$. Fix $k \in K$ and write $B_k = B_{k,j_k}$. By
Corollary~\ref{cor:step},
$(B_k + \sigma_k I)(x^{k+1} - x^k)
= \sum_{i=1}^n (\nu^\ell_i - \nu^u_i) e_i
- \sum_{i \in I_\delta(x^k)} \mu_i \nabla f_i(x^k)$, hence
\[
\big\| (B_k + \sigma_k I)(x^{k+1} - x^k) \big\|
= \Big\| \sum_{i \in I_\delta(x^k)} \mu_i \nabla f_i(x^k)
+ \sum_{i=1}^n (\nu^u_i - \nu^\ell_i) e_i \Big\| > \epsilon,
\]
the inequality holding because $C_{\delta,\epsilon}$ fails at $x^{k+1}$. Since
$\|B_k\| \leq M$,
\[
\epsilon < \| (B_k + \sigma_k I)(x^{k+1} - x^k) \|
\leq (M + \sigma_k) \|x^{k+1} - x^k\|,
\]
and with \eqref{cota-sigma},
$\|x^{k+1} - x^k\| > \epsilon / \big( M + \gamma\max\{L + 2\alpha, 9c_x^2/(2\delta)\} \big)$.
As $x^{k+1}$ satisfies \eqref{ovo-armijo},
\[
f(x^{k+1}) \leq f(x^k) - \alpha\|x^{k+1} - x^k\|^2
< f(x^k) - \frac{\alpha\epsilon^2}
{\big( M + \gamma\max\{L + 2\alpha, 9c_x^2/(2\delta)\} \big)^2}.
\]
Summing over $k \in K$ and using $f(x) \geq \flow$ gives
$f(x^0) - \flow > |K| \cdot \alpha\epsilon^2 /
\big( M + \gamma\max\{L + 2\alpha, 9c_x^2/(2\delta)\} \big)^2$, from which
\eqref{cota-iter} follows.
\end{proof}

\section{Numerical experiments} \label{sec:numerical}

We illustrate the method on parameter-estimation problems in which the data
contain an unknown number of outliers, following the OVO methodology of
\cite{alvarezovo}. Given a model $y(t,x)$ with parameters $x\in\Omega$ and data
$(t_i,y_i)$, $i=1,\dots,m$, we set
\[
f_i(x)=\tfrac12\big(y(t_i,x)-y_i\big)^2,\qquad i=1,\dots,m.
\]
For a prescribed number $o$ of outliers, solving \eqref{ovo-problem} with
$p=m-o$ fits the parameters to the $p$ best-explained observations while
discarding the $o$ worst. Since the true number of outliers is unknown, we
solve a sequence of OVO problems with increasing $o$; a sharp drop in $f(x^*)$
identifies the number of outliers present.

Our aim is to assess the effect of the curvature matrix $B$. We therefore
compare two instances of Algorithm~\ref{sec:method}.1 that differ \emph{only}
in $B$: the first-order method ($B_{k,j}=0$, which coincides with
\cite{alvarezovo}) and the quasi-Newton variant, in which $B_{k,j}$ is built
from the Gauss--Newton curvature of the active component. For the least-squares
functions above, $\nabla f_i(x)=r_i(x)\nabla r_i(x)$ with
$r_i(x)=y(t_i,x)-y_i$, so the Gauss--Newton curvature
$\nabla r_i(x)\nabla r_i(x)^T\succeq0$ is available from quantities already
computed for the model and requires no second derivatives. To respect the
bound $\|B_{k,j}\|\le M$ of Step~2, we scale the curvature so that its spectral
norm does not exceed $M$; like $\delta$, the value of $M$ is chosen by trial
and error.

As in \cite{alvarezovo}, subproblem \eqref{ovo-subpro} is reformulated as the
smooth nonlinear program
\[
\minimize_{x\in\Omega,\,w}\ w \quad\text{s.t.}\quad
\nabla f_i(x^k)^T(x-x^k)
+\tfrac12\big[(x-x^k)^TB_{k,j}(x-x^k)+\sigma_{k,j}\|x-x^k\|^2\big]\le w,
\quad i\in I_\delta(x^k),
\]
solved with Algencan \cite{bmbook,bmcomper}, whose multipliers furnish
$\mu,\nu^\ell,\nu^u$ for checking $C_{\delta,\epsilon}$. Throughout we use
$\sigma_{\min}=0.1$, $\alpha=10^{-8}$, $\gamma=5$, and $\epsilon=10^{-4}$. Both
methods are run from identical starting points with the same random seed, so
that any difference is attributable to $B$.

The codes were implemented in Fortran 90. Tests were conducted on a computer
with a 13th-generation Intel Core i7-13700KF processor (up to $5.4$\,GHz) and
$32$\,GB $4800$\,MT/s DDR5 RAM,
running Linux (Arch Linux, kernel 7.0.10). Code was compiled by the GNU Fortran
compiler (version 16.1.1) with the \texttt{-O3} optimization directive enabled.

\subsection{A model linear in the parameters} \label{subsec:poly}

Following \cite{dunder1}, we fit the third-degree polynomial
$y(t,x)=x_1+x_2t+x_3t^2+x_4t^3$, with $-10\le x_i\le10$, to a set of $m=46$
observations containing $10$ outliers, taken from \cite{dunder1}. The
least-squares fit to the contaminated data is
$\bar x\approx(6.46,\,2.71,\,-7.54,\,2.16)^T$. As the OVO problem has many
non-global minimizers, for each $o\in\{0,\dots,12\}$ we run the algorithm from
$100$ starting points $x^0_i=\bar x_i+r\,|\bar x_i|$, with $r\sim U[-0.5,0.5]$
(projected onto $\Omega$), and report the best $x^*$. We take $\delta=0.1$ and
$M=1$.

Table~\ref{tab:poly} reports, for each $o$ and for both methods, the best value
$f(x^*)$ over the $100$ starting points, the iterations and evaluations of the
single run that attained it, and the mean and median iterations and evaluations
over all $100$ runs. Both methods identify the $10$ outliers
through the sharp drop of $f(x^*)$ between $o=9$ and $o=10$, and return solutions
of the same quality. The ``best run'' counts are a single, point-initialisation
dependent sample and are reported only to show at what cost that particular
solution was reached; the fair cost comparison is the mean and median over all
$100$ runs, on which the two methods are essentially identical. Summing over $o$
the per-$o$ means reported in
Table~\ref{tab:poly} gives $368$ iterations and $1836$ evaluations for the
first-order method against $369$ and $1858$ for the quasi-Newton variant; the
corresponding sums of the per-$o$ medians (the median columns of the table)
agree just as closely, $350$ and $1760$ against $352$ and $1792$. The two methods thus
differ only within the run-to-run variability of the multistart
procedure. This is expected: on this affine model the Gauss--Newton curvature
coincides with the exact, constant Hessian, so the explicit curvature term is
redundant and $B$ has no discernible effect.

\begin{table}[ht!]
\centering

\resizebox{\textwidth}{!}{%
\begin{tabular}{c ccc cc cc ccc cc cc}
\hline
 & \multicolumn{7}{c}{First-order ($B=0$)} & \multicolumn{7}{c}{Quasi-Newton ($M=1$)} \\
\cline{2-8}\cline{9-15}
 & \multicolumn{3}{c}{best run} & \multicolumn{2}{c}{\#it} & \multicolumn{2}{c}{\#fcnt}
 & \multicolumn{3}{c}{best run} & \multicolumn{2}{c}{\#it} & \multicolumn{2}{c}{\#fcnt} \\
\cline{2-4}\cline{5-6}\cline{7-8}\cline{9-11}\cline{12-13}\cline{14-15}
$o$ & $f(x^*)$ & \#it & \#fcnt & mean & med & mean & med
    & $f(x^*)$ & \#it & \#fcnt & mean & med & mean & med \\
\hline
0  & 1.363e+01 & 20 & 120 & 31 & 30.0 & 151 & 149.5 & 1.362e+01 & 45 & 205 & 30 & 29.5 & 147 & 143.5 \\
1  & 1.145e+01 & 45 & 238 & 30 & 30.0 & 146 & 142.0 & 1.145e+01 & 20 &  77 & 31 & 30.0 & 155 & 152.5 \\
2  & 1.004e+01 & 22 & 114 & 31 & 28.5 & 155 & 152.5 & 1.005e+01 & 17 &  84 & 30 & 28.5 & 152 & 150.0 \\
3  & 9.535e+00 & 40 & 182 & 29 & 28.0 & 147 & 143.5 & 9.530e+00 & 15 &  84 & 28 & 27.0 & 143 & 137.0 \\
4  & 9.013e+00 & 64 & 341 & 29 & 27.5 & 140 & 132.5 & 9.014e+00 & 33 & 170 & 28 & 27.0 & 138 & 132.0 \\
5  & 8.455e+00 & 38 & 170 & 29 & 27.0 & 144 & 133.0 & 8.460e+00 & 32 & 158 & 30 & 29.0 & 149 & 143.0 \\
6  & 7.436e+00 & 38 & 199 & 28 & 24.0 & 143 & 126.0 & 7.434e+00 & 40 & 193 & 29 & 26.5 & 147 & 135.5 \\
7  & 6.921e+00 & 31 & 151 & 30 & 30.0 & 151 & 150.5 & 6.907e+00 & 19 &  89 & 31 & 31.0 & 161 & 155.0 \\
8  & 5.503e+00 & 39 & 158 & 28 & 27.0 & 139 & 139.0 & 5.501e+00 & 32 & 139 & 28 & 26.0 & 143 & 140.5 \\
9  & 4.176e+00 & 33 & 150 & 26 & 24.0 & 132 & 117.5 & 4.190e+00 & 55 & 282 & 25 & 23.5 & 126 & 119.5 \\
10 & 3.120e-02 & 29 & 170 & 26 & 23.5 & 130 & 122.0 & 4.077e-02 & 35 & 182 & 26 & 24.0 & 133 & 123.0 \\
11 & 3.452e-02 & 27 & 137 & 25 & 25.0 & 126 & 122.5 & 2.275e-02 & 21 &  89 & 26 & 24.5 & 129 & 124.5 \\
12 & 3.228e-02 & 23 & 142 & 26 & 25.5 & 132 & 130.0 & 3.421e-02 & 22 & 116 & 27 & 26.0 & 135 & 136.0 \\
\hline
$\Sigma$ & & 449 & 2272 & 368 & 350.0 & 1836 & 1760.5 & & 386 & 1868 & 369 & 352.5 & 1858 & 1792.0 \\
\hline
\end{tabular}%
}
\caption{Third-degree polynomial ($m=46$, $\delta=0.1$): first-order method
($B=0$) versus the quasi-Newton variant ($M=1$). For each $o$, the ``best run''
block reports the best $f(x^*)$ over $100$ starting points together with the
outer iterations (\#it) and function evaluations (\#fcnt) of the single run that
attained it; the remaining columns give the mean and median \#it and \#fcnt over
all $100$ starting points. The $\Sigma$ row sums each count over $o$.}
\label{tab:poly}
\end{table}

The absence of any gain from curvature on this problem is expected, and
instructive. The model $y(t,x)$ is affine in the parameters $x$, so each
residual $r_i$ is affine, $\nabla r_i$ is constant, and the component Hessians
$\nabla^2 f_i=\nabla r_i\nabla r_i^T$ are constant and exact: here the
Gauss--Newton curvature is the true curvature. Each $f_i$ is thus a convex
quadratic, for which the adaptive regularization $\sigma_{k,j}$ of the
first-order model already produces well-scaled steps; the explicit curvature is
largely redundant and merely reshapes the step without reducing the iteration
count. Moreover, the dominant difficulty of \eqref{ovo-problem} on this
instance is not smooth nonconvexity --which curvature would mitigate-- but the
combinatorial selection of which observations to discard, which produces the
many local minima reflected in the multistart strategy and is unaffected by
$B$. This is consistent with the analysis of Section~\ref{sec:complexity},
where $B$ enters only the constant, and not the order, of the complexity
bounds: on a problem of this type curvature can be expected to neither help nor
harm. We turn next to a model that is genuinely nonlinear in the parameters, where
the first-order model is a poorer local approximation, to examine whether the
additional curvature translates into a practical advantage.

\subsection{A model nonlinear in the parameters} \label{subsec:bard}

Following \cite{mgh}, we fit the Bard function, whose $i$th residual is
\begin{equation} \label{bard-model}
r_i(x) = x_1 + \frac{u_i}{v_i x_2 + w_i x_3} - y_i, \qquad
u_i = t_i,\ \ v_i = 16 - t_i,\ \ w_i = \min(u_i, v_i),
\end{equation}
where $t_i$ is the abscissa of the $i$th observation, and set
$f_i(x) = \tfrac12 r_i(x)^2$; the model has $n = 3$ parameters. The original
data set of \cite{mgh} consists of $m_0 = 15$ observations at $t_i = i$,
$i = 1, \dots, 15$. The model is nonlinear in $x_2$ and $x_3$, which
enter through the denominator, so the component Hessians
$\nabla^2 f_i = \nabla r_i \nabla r_i^T + r_i \nabla^2 r_i$ are no longer
constant and the Gauss--Newton term carries genuine curvature information. We
take the box $x_1 \in [-10, 10]$, $x_2, x_3 \in [0.1, 10]$; the lower bounds on
$x_2, x_3$ keep the denominator in \eqref{bard-model} positive and exclude the
spurious minimizer of the Bard function at $(0.8406\ldots, -\infty, -\infty)$.

To the $m_0 = 15$ clean observations we add four outliers at random abscissas
$t_i \in (2, 14)$, two lying above and two below the model curve, with random
magnitudes clearly separated from the clean-data noise; drawn with a fixed
random seed, they yield $m = 19$ observations. As in Section~\ref{subsec:poly}
we run the two variants from $100$ starting points obtained by perturbing the
standard initial guess $x^0 = (1,1,1)$, and report the best $x^*$; we take
$\delta = 0.1$ and $M = 1$. For each $o \in \{0, \dots, 6\}$ we solve
\eqref{ovo-problem} with $p = m - o$.

Table~\ref{tab:bard} reports the results. As the true number of outliers is
unknown, one solves a sequence of problems with increasing $o$ and looks for a
sharp drop in $f(x^*)$; here $f(x^*)$ decreases only moderately for $o \le 3$
and then drops by more than two orders of magnitude between $o = 3$ and $o = 4$,
after which it remains negligible (Figure~\ref{fig:bard}(b)). Both methods exhibit the same drop at $o = 4$
and both identify exactly the four added outliers. For $o \le 3$ they return
essentially identical values of $f(x^*)$; for $o \ge 4$, once the outliers are
removed, both reduce $f(x^*)$ to the residual floor of the clean data, the
first-order method reaching marginally lower values there, within the
variability of the multistart procedure. Figure~\ref{fig:bard}(a) shows the fit
obtained at $o = 4$: the model follows the fifteen clean observations while the
four outliers, marked with circles, are discarded.

\begin{table}[ht!]
\centering
\resizebox{\textwidth}{!}{%
\begin{tabular}{c ccc cc cc ccc cc cc}
\hline
 & \multicolumn{7}{c}{First-order ($B=0$)} & \multicolumn{7}{c}{Quasi-Newton ($M=1$)} \\
\cline{2-8}\cline{9-15}
 & \multicolumn{3}{c}{best run} & \multicolumn{2}{c}{\#it} & \multicolumn{2}{c}{\#fcnt}
 & \multicolumn{3}{c}{best run} & \multicolumn{2}{c}{\#it} & \multicolumn{2}{c}{\#fcnt} \\
\cline{2-4}\cline{5-6}\cline{7-8}\cline{9-11}\cline{12-13}\cline{14-15}
$o$ & $f(x^*)$ & \#it & \#fcnt & mean & med & mean & med
    & $f(x^*)$ & \#it & \#fcnt & mean & med & mean & med \\
\hline
0 & 2.643e+00 & 367 & 394 & 361 & 362.5 & 382 & 385.0 & 2.643e+00 & 393 & 411 & 399 & 397.5 & 416 & 414.5 \\
1 & 1.879e+00 &   8 &  30 & 152 &  10.5 & 166 &  29.5 & 1.876e+00 &  16 &  45 &  39 &   7.0 &  54 &  22.0 \\
2 & 6.765e-01 &   5 &  17 &  32 &   7.0 &  45 &  23.0 & 6.693e-01 &   4 &  12 &  38 &   7.0 &  49 &  18.0 \\
3 & 3.511e-01 &   7 &  21 &  13 &   8.0 &  28 &  22.0 & 3.549e-01 &   8 &  27 &  11 &   6.0 &  20 &  15.0 \\
4 & 2.086e-03 &   7 &  24 &  10 &   8.0 &  19 &  17.5 & 2.785e-03 &   2 &   5 &   8 &   6.0 &  13 &  11.0 \\
5 & 6.578e-05 &   6 &  27 &  16 &   5.0 &  22 &  11.0 & 1.021e-04 &   6 &   9 &   8 &   6.0 &  11 &   8.0 \\
6 & 6.319e-05 &  13 &  21 &   9 &   5.0 &  16 &  12.0 & 9.543e-04 &   7 &   8 &  10 &   6.0 &  13 &   9.0 \\
\hline
$\Sigma$ & & 413 & 534 & 593 & 406.0 & 678 & 500.0 & & 436 & 517 & 513 & 435.5 & 576 & 497.5 \\
\hline
\end{tabular}%
}
\caption{Bard function ($m=19$, $\delta=0.1$) with four added outliers,
two above and two below the model curve: first-order method ($B=0$) versus the
quasi-Newton variant ($M=1$). For each $o$, the ``best run'' block reports the
best $f(x^*)$ over $100$ starting points together with the outer iterations
(\#it) and function evaluations (\#fcnt) of the single run that attained it; the
remaining columns give the mean and median \#it and \#fcnt over all $100$
starting points. The $\Sigma$ row sums each count over $o$.}
\label{tab:bard}
\end{table}

\begin{figure}[ht!]
\centering
\includegraphics[width=\textwidth]{bard_panels.pdf}
\caption{Bard function with four added outliers, two above and two below the
model curve. (a) The fit returned at $o=4$ (solid line) follows the fifteen
clean observations (dots), while the four points identified as outliers
(circles) are discarded. (b) Best $f(x^*)$ over the $100$ starting points as a
function of the number $o$ of discarded observations (logarithmic vertical axis;
the two variants coincide): $f(x^*)$ decreases only moderately up to $o=3$ and
drops by more than two orders of magnitude at $o=4$, after which it stays at the
residual floor of the clean data, the onset of which identifies the four
outliers.}
\label{fig:bard}
\end{figure}

Unlike the polynomial instance, the counts now differ between the two methods,
and the difference has structure. The aggregate over all $o$ is dominated by
$o = 0$---the minimax over all $m = 19$ observations, outliers included---a hard
nonconvex problem where both methods are expensive (about $360$--$400$ mean
iterations) and the first-order method is in fact slightly cheaper. Once at
least the outliers are being discarded ($o \ge 1$), the regime for which the
method is intended, the quasi-Newton variant is consistently more efficient.
Summing the per-$o$ means of Table~\ref{tab:bard} over $o = 1, \dots, 6$ gives
$232$ iterations and $296$ evaluations for the first-order method against $114$
and $160$ for the quasi-Newton variant; the corresponding sums of the per-$o$
medians are $43.5$ and $115$ against $38$ and $83$. The mean gap is large---the
first-order method uses about twice as many iterations and evaluations---but is
amplified by a heavy tail of slow first-order runs on the hardest instances (at
$o = 1$ the mean iteration count is $152$ for the first-order method against
$39$ for the quasi-Newton variant), which the curvature mitigates. The median,
robust to this tail, still favours the quasi-Newton variant, though more
modestly. Either way, in this regime the curvature lowers the cost.

Curvature thus helps precisely where the first-order model is a poor local
approximation---a nonlinear model with outliers still present---while it neither
helps on the affine model of Section~\ref{subsec:poly} nor on the full-data
minimax at $o = 0$. This is consistent with the analysis of
Section~\ref{sec:complexity}, where $B$ affects only the constants of the
complexity bounds: the gain is a constant-factor one, realized here in the
outlier-rich regime, not a change in the order. Two caveats temper it. The
evaluation count does not reflect the extra per-iteration work of assembling and
spectrally processing the curvature matrix, so the wall-clock benefit is smaller
than the counts suggest; and once the outliers are removed ($o \ge 4$) both
methods reach the residual floor of the clean data, where the first-order
variant attains marginally lower $f(x^*)$ values, within the multistart
variability.

\section{Final remarks} \label{sec:final}

We have extended the first-order regularized method of \cite{alvarezovo} for
box-constrained order-value optimization by enriching its model with
quasi-Newton curvature. The resulting method is well defined and, under the same
assumptions, attains the same worst-case iteration- and evaluation-complexity
bounds as its first-order counterpart, of which it is a strict generalization.
The curvature is thus accommodated at no cost in the theoretical guarantees, and
standard quasi-Newton updates such as BFGS or SR1 fit directly into the
framework.

The numerical experiments delineate where the curvature helps and where it does
not. On the affine model the Gauss--Newton curvature is exact but redundant, and
the two variants are indistinguishable. On the nonlinear model the picture
depends on the regime: for the full-data minimax the curvature does not help,
but once outliers are being discarded---the regime for which order-value
optimization is intended---the quasi-Newton variant is consistently more
efficient, reducing both the mean and the median iteration and evaluation counts
and, in particular, curbing a heavy tail of slow first-order runs on the hardest
instances. The advantage is a constant-factor one, realized where the
first-order model is a poor local approximation, and is consistent with the
complexity analysis, in which $B$ enters only the constants and not the order of
the bounds. It should be read with two caveats: the per-iteration cost of
assembling and spectrally processing the curvature matrix is not reflected in
the evaluation count, so the wall-clock benefit is smaller than the counts
suggest; and the solution quality is essentially the same for both variants.

The reason the curvature does not change the order is structural. The Tikhonov
term $\sigma\|x - \bar x\|^2$ makes the model an upper approximation of $f$ once
$\sigma$ is of order $L$, which yields a step of size
$\|x^{k+1} - x^k\| = O(\epsilon)$ and a per-iteration decrease of order
$\epsilon^2$, hence the $O(\epsilon^{-2})$ iteration bound. A natural way to
genuinely improve the order is to replace the quadratic regularization by a
cubic one, $\tfrac{\sigma}{3}\|x - \bar x\|^3$, in the spirit of the
cubic-regularization method of \cite{nesterov} and the high-order regularized
models of \cite{bgmst1, bgms}. With accurate second-order information---the true
or Gauss--Newton Hessian of each active component rather than an arbitrary
bounded $B$---and under the additional assumption that the component Hessians are
Lipschitz continuous, the cubic model is an upper approximation with a step of
size $O(\epsilon^{1/2})$ and a per-iteration decrease of order $\epsilon^{3/2}$,
which would improve the iteration bound to $O(\epsilon^{-3/2})$. Realizing this
for the order-value function is not a direct transfer: the objective is
nonsmooth, the complexity carries the two parameters $\delta$ and $\epsilon$, and
the approximate optimality condition $C_{\delta,\epsilon}$ is stated in terms of
active gradients. Extending the analysis to a cubically regularized model,
tracking the interaction of the improved $\epsilon$-dependence with $\delta$, is
a promising direction in which the curvature would acquire a provable role, at
the price of a stronger smoothness assumption, accurate second-order information,
and a cubically regularized subproblem.

\clearpage

\begin{thebibliography}{99}

\bibitem{dunder1} R. Andreani, C. Dunder and J. M. Mart\'{\i}nez, Order-Value Optimization: Formulation and solution by means of a primal Cauchy method, \textit{Mathematical Methods of Operations Research} 58, pp. 387--399, 2003.

\bibitem{dunder2} R. Andreani, C. Dunder and J. M. Mart\'{\i}nez, Nonlinear-programming reformulation of the order-value optimization problem, \textit{Mathematical Methods of Operations Research} 61, pp. 365--384, 2005.

\bibitem{yano1} R. Andreani, J. M. Mart\'{\i}nez, M. Salvatierra, and F. S. Yano, Quasi-Newton methods for Order-Value Optimization and Value-at-Risk calculations, \textit{Pacific Journal of Optimization} 2, pp. 11--33, 2006.

\bibitem{yano2} R. Andreani, J. M. Mart\'{\i}nez, L. Mart\'{\i}nez, and F. S. Yano, Low order-value optimization and applications, \textit{Journal of Global Optimization} 43, pp. 1--22, 2009.

\bibitem{salva} R. Andreani, J. M. Mart\'{\i}nez, M. Salvatierra, and F. S. Yano, Global order-value optimization by means of a multistart harmonic oscillator tunneling strategy, in \textit{Global Optimization: From Theory to Implementation}, L. Liberti and N. Maculan (eds.), Springer, Boston, MA, 2006, pp. 379--404.

\bibitem{bgms} E. G. Birgin, J. L. Gardenghi, J. M. Mart\'{\i}nez, and S. A. Santos, On the use of third-order models with fourth-order regularization for unconstrained optimization, \textit{Optimization Letters} 14, pp. 815--838, 2020.

\bibitem{bgmst2} E. G. Birgin, J. L. Gardenghi, J. M. Mart\'{\i}nez, S. A. Santos, and Ph. L. Toint, Evaluation complexity for nonlinear constrained optimization using unscaled KKT conditions and high-order models, \textit{SIAM Journal on Optimization} 26, pp. 951--967, 2016.

\bibitem{bgmst1} E. G. Birgin, J. L. Gardenghi, J. M. Mart\'{\i}nez, S. A. Santos, and Ph. L. Toint, Worst-case evaluation complexity for unconstrained nonlinear optimization using high-order regularized models, \textit{Mathematical Programming} 163, pp. 359--368, 2017.

\bibitem{bmbook} E. G. Birgin and J. M. Mart\'{\i}nez, \textit{Practical Augmented Lagrangian Methods for Constrained Optimization}, Society for Industrial and Applied Mathematics, Philadelphia, 2014.

\bibitem{bmcomper} E. G. Birgin and J. M. Mart\'{\i}nez, Complexity and performance of an Augmented Lagrangian algorithm, \textit{Optimization Methods and Software} 35, pp. 885--920, 2020.

\bibitem{cogtbook} C. Cartis, N. I. M. Gould, and Ph. L. Toint, \textit{Evaluation Complexity of Algorithms for Nonconvex Optimization: Theory, Computation, and Perspectives}, SIAM, Philadelphia, PA, 2022.

\bibitem{gauvin} J. Gauvin, A necessary and sufficient regularity condition to have bounded multipliers in nonconvex programming, \textit{Mathematical Programming} 12, pp. 136--138, 1977.
  
\bibitem{jorion} P. Jorion, \textit{Value at {R}isk: {T}he new benchmark for managing financial risk}, 3rd. ed., McGraw-Hill, 2009.

\bibitem{mtop} J. M. Mart\'{\i}nez, Generalized Order-Value Optimization, \textit{TOP} 20, pp. 75--98, 2012.

\bibitem{mgh} J. J. Mor\'e, B. S. Garbow, and K. E. Hillstrom, Testing unconstrained optimization software, \textit{ACM Transactions on Mathematical Software} 7, pp. 17--41, 1981.
  
\bibitem{nesterov} Y. Nesterov and B. T. Polyak, Cubic regularization of Newton method and its global performance, \textit{Mathematical Programming} 108, 177--205, 2006.

\bibitem{jiang} Z. Jiang, Q. Hu, and X. Zheng, Optimality condition and complexity of order-value optimization problems and low order-value optimization problems, \textit{Journal of Global Optimization} 69, pp. 511--523, 2017.

\bibitem{alvarezovo} G. Q. Álvarez and E. G. Birgin, A first-order regularized approach to the order-value optimization problem, \textit{Optimization Methods and Software} 40, pp. 650--674, 2025.
    
\end{thebibliography}

\end{document}